\section*{Chapter \ref{ch:iv1:timed-test-3} --- Timed Test 3}
\markboth{\protect\omTocMark{}{Answers: Timed Test 3}}{\protect\omTocMark{}{Answers: Timed Test 3}}

\begin{solution}{iq:iv1:timed-test-3:1}
The average is $7.3$ boxes.

Let $T_i$ be the wait for toy $i$. The time to get all is $\max T_i$, and inclusion--exclusion\index{inclusion--exclusion} over minimums gives
\[
\mathbb{E}\max = \sum_i\frac{1}{p_i} - \sum_{i < j}\frac{1}{p_i + p_j} + \frac{1}{p_1 + p_2 + p_3}.
\]
The first sum is $2 + 3 + 6 = 11$. The pairs have $\frac56$, $\frac23$ and $\frac12$, which give $1.2 + 1.5 + 2 = 4.7$. The last term is 1. So the mean is $11 - 4.7 + 1 = 7.3$.

Sanity check: it must exceed 6, the wait for the rarest toy, and it does so by a little. Equal probabilities would give $5.5$, so uneven ones cost more.

\iqwrong{Using the equal-probability formula $3(1 + \frac12 + \frac13)$, which gives $5.5$.}
\end{solution}

\begin{solution}{iq:iv1:timed-test-3:2}
A total of 10 is more likely: $\frac{27}{216}$ against $\frac{25}{216}$.

Both totals have six unordered partitions, which is why they can look equal. But the partitions have different numbers of orderings\index{counting ordered outcomes}.

For 9: 126 (6), 135 (6), 144 (3), 225 (3), 234 (6), 333 (1), a total of 25.

For 10: 136 (6), 145 (6), 226 (3), 235 (6), 244 (3), 334 (3), a total of 27.

A check: by the symmetry $x \to 7 - x$, totals 10 and 11 have equal probability, as do 9 and 12. And an enumeration over the 216 outcomes confirms both numbers.

\iqwrong{Counting partitions as equally likely, which treats 333 like 126. The outcomes are the 216 ordered triples, and only they are equally likely.}
\end{solution}

\enlargethispage{-3\baselineskip}%
\begin{solution}{iq:iv1:timed-test-3:3}
Both are $1.75$ bits.

Huffman's construction\index{Huffman code} merges the two rarest symbols and repeats. It gives code lengths 1, 2, 3 and 3. The expected length is $\frac12 \cdot 1 + \frac14 \cdot 2 + \frac18 \cdot 3 + \frac18 \cdot 3 = 1.75$.

The entropy\index{entropy} is $\sum p\log_2(1/p)$, with the same numbers, since each $\log_2(1/p)$ is an integer. So the code achieves the theoretical floor exactly.

Compare a fixed-length code, which takes 2 bits. The gain comes from using short words for frequent symbols.

\iqfollow{When does the code not reach the entropy?}{When the probabilities are not powers of $\frac12$. For a source with probabilities $0.9$ and $0.1$, the code takes 1 bit a symbol, while the entropy is $0.47$. Coding blocks of symbols closes this gap.}
\end{solution}

\begin{solution}{iq:iv1:timed-test-3:4}
The probability is $\frac12(1 - \ln 2) \approx 0.153$.

Condition on $X = x$\index{total probability}. For $x \le \frac12$, the event is impossible. For $x > \frac12$, it has chance $(x - \frac12)/x = 1 - \frac{1}{2x}$. So
\[
P = \int_{1/2}^1\left(1 - \frac{1}{2x}\right)dx = \frac12 - \frac12\ln 2.
\]

A second route uses the density of $Y$, which is $-\ln y$ on $(0, 1)$, found by integrating $1/x$ from $y$ to 1. Then $P = \int_{1/2}^1(-\ln y)\,dy = \frac12 - \frac12\ln 2$, the same.

The mean of $Y$ is $\frac14$, so a chance of $0.15$ to exceed $\frac12$ is reasonable.

\iqwrong{Using $X$'s mean, $\frac12$, as if $Y$ were uniform on $[0, \frac12]$, which gives 0. A value cannot replace a distribution here.}
\end{solution}

\begin{solution}{iq:iv1:timed-test-3:5}
A wins with probability $\frac38$.

Each full round, A wins with chance $\frac16$. If A fails, B wins with chance $\frac13$. So the game continues to another round with chance $\frac56 \cdot \frac23 = \frac59$\index{geometric series}.

Let $a$ be A's chance. Then $a = \frac16 + \frac59a$, which gives $a = \frac16 \cdot \frac94 = \frac38$.

A check: B's chance is $\frac56 \cdot \frac13/(\frac49) = \frac58$, and the two add to 1. So going first is not enough to offset a target that is twice as easy. With equal targets of a 6, A would have $\frac{6}{11}$.

\iqwrong{Adding up only A's first roll, or giving $\frac13$ by comparing the two chances directly, $\frac16/(\frac16 + \frac13)$. That ignores who rolls first.}
\end{solution}

\begin{solution}{iq:iv1:timed-test-3:6}
A right triangle has probability $\frac35$, and an equilateral one has probability $\frac{1}{10}$.

There are $\binom63 = 20$ triangles. The hexagon's vertices lie on a circle, so a triangle has a right angle exactly when two of its vertices form a diameter\index{Thales' theorem}. There are 3 diameters. Each can be combined with any of the 4 other vertices, which gives 12 triangles, all different, since a triangle cannot contain two diameters.

So the right triangles have $\frac{12}{20}$. The equilateral triangles are the two formed by alternate vertices, so $\frac{2}{20}$. The remaining 6 are isosceles with an obtuse angle: three neighbouring vertices. An enumeration confirms the counts.

\iqwrong{Trying to compute angles for each case. The diameter criterion gives the count at once.}
\end{solution}

\begin{solution}{iq:iv1:timed-test-3:7}
The probability is $\frac18$, and the expected number of cycles is $H_8 = \frac{761}{280} \approx 2.72$.

Follow the cycle from 8. It closes at the first step with chance $\frac18$, the value $\pi(8) = 8$. Otherwise it continues to another element, and closes at the next step with chance $\frac17$, and so on. The chance that it has length 3 is $\frac78 \cdot \frac67 \cdot \frac16 = \frac18$\index{cycle length}. All lengths from 1 to 8 are equally likely, by the same cancellation.

For the number of cycles, use that each element $i$ in a cycle of length $L_i$ contributes $1/L_i$, with total 1 per cycle. So the count is $\sum_i1/L_i$, and its mean is $8 \cdot \frac18(1 + \frac12 + \dots + \frac18) = H_8$. Counting by brute force over all permutations gives the same.

\iqfollow{What is the mean length of the cycle holding a given element?}{$(n + 1)/2 = 4.5$, the mean of the uniform on 1 to 8. Note that this is larger than the average cycle length, $8/H_8 = 2.94$, an inspection effect.}
\end{solution}

\begin{solution}{iq:iv1:timed-test-3:8}
About $0.196$ tests per person, a saving of about 80\%.

The pool costs one test, or $\frac{1}{10}$ per person. With chance $1 - 0.99^{10}$, the pool is positive, and each person takes one more test\index{group testing}. So the mean per person is
\[
\frac{1}{10} + 1 - 0.99^{10} = 0.1 + 0.0956 = 0.1956.
\]
For the value, $0.99^{10} \approx e^{-0.1} \approx 0.904$.

The best pool size balances the two terms: for $f(k) = 1/k + 1 - 0.99^k$, the minimum is near $k \approx 1/\sqrt{p} = 10$. An exact search shows that pools of 10 and 11 are almost the same. So 10 is a good design.

\iqfollow{If the rate rises to 30\%, is pooling still useful?}{Barely. A pool of 2 costs $1.01$ per person, the best pool, of 3, costs $0.99$, and larger pools cost more. So it no longer pays to bother.}
\end{solution}

\begin{solution}{iq:iv1:timed-test-3:9}
The probability is $\frac16(\frac76)^5 \approx 0.360$, and for large $n$ it tends to $\frac27$.

Let $p_n$ be the chance of hitting $n$, with $p_0 = 1$. Conditioning on the last roll\index{renewal equation}, for $n \le 6$,
\[
p_n = \frac16(p_{n - 1} + \dots + p_0).
\]
So $p_1 = \frac16$, and $p_{n + 1} = p_n + \frac16p_n = \frac76p_n$ for $n \ge 1$ up to 5. Thus $p_6 = \frac16(\frac76)^5$. An exact recursion in fractions gives this number.

For large $n$, renewal theory says that hits occur at a rate of one per mean step, and the mean step is $3.5$. So $p_n \to 1/3.5 = \frac27 = 0.286$.

Note that $p_6$ is the largest of the $p_n$: from 6 onwards, the bound on the number of earlier terms in the recursion takes effect.

\iqwrong{Enumerating the compositions of 6 and weighting each by $(\frac16)^k$. That works, but it is the slow form of the same answer.}
\end{solution}

\begin{solution}{iq:iv1:timed-test-3:10}
The average area is $\frac16$, and the area exceeds $\frac{3}{16}$ with chance $\frac12$.

With the break at $x$, uniform on $[0, 1]$, the area is $x(1 - x)$, whose mean is $\frac12 - \frac13 = \frac16$. The area exceeds $\frac{3}{16}$ when $(x - \frac14)(x - \frac34) < 0$, that is when $x$ lies between $\frac14$ and $\frac34$, and this has chance $\frac12$.

A check: the shorter piece $S$ has mean $\frac14$ and the longer $L$ has mean $\frac34$, with product $\frac{3}{16}$. But $L = 1 - S$, so their covariance\index{covariance} is $-\operatorname{Var}(S) = -\frac{1}{48}$, and $\frac{3}{16} - \frac{1}{48} = \frac16$.

\iqwrong{Giving $\frac{3}{16}$, the product of the mean pieces, or $\frac14$, the area at the mean break. The first ignores that the pieces are negatively correlated, and the second is the largest area possible, reached by a single break point.}

\iqwhy{The mean of a product is the product of the means only without correlation. The average of price times volume is not the average price times the average volume.}
\end{solution}

\begin{solution}{iq:iv1:timed-test-3:11}
The straddle costs about $6.4$, and the daily break-even move is about 1, or 1\%.

The call is worth $0.4 \times 100 \times 0.16 \times 0.5 = 3.2$. At the money with zero rates, put and call have the same value, by parity. So the straddle costs $6.4$.

The factor comes from $2\phi(0)\frac12 \approx 0.4$\index{at-the-money approximation}: the call is $S(2\Phi(\frac12\sigma\sqrt{T}) - 1)$, and it is linear in a small argument.

A delta-hedged option earns $\frac12\Gamma S^2(r^2 - \sigma^2\Delta t)$ per step, so it breaks even when the daily return equals $\sigma\sqrt{\Delta t} = 16\%/\sqrt{252} = 1.01\%$\index{gamma}. That is the rule of 16: annual vol over 16 is the daily move.

\iqwrong{Comparing the $\pm 6.4$ break-even at expiry with daily moves. A hedged holder's result depends on the sum of squared moves, not on where the stock ends.}

\iqfollow{What does the straddle cost at one year?}{$6.4 \times 2 = 12.8$, since the time scale doubles the root.}
\end{solution}

\begin{solution}{iq:iv1:timed-test-3:12}
The average is $\frac{11}{3}$.

Use indicators\index{indicator variables}. An interior position, with two neighbours, has the largest of three random distinct values with chance $\frac13$. An end position, with one neighbour, has chance $\frac12$.

There are 8 interior positions and 2 ends, so the mean is $\frac83 + 1 = \frac{11}{3}$. Brute force for small $n$ follows the general formula $(n + 1)/3$.

The result is about one in three positions, which is the rate at which a random series makes a peak. So if a chart seems to have fewer turning points than this, the series has a trend or a memory.

\iqfollow{What is the count of runs of ascents, for $n = 10$?}{The number of ascents has mean $4.5$, and maximal ascending runs number one more than the descents, so $5.5$ on average. These are the same indicators with pairs in place of triples.}
\end{solution}
